\section{Detailed Measurements}
\label{app:meas}

This appendix expands the headline of \Cref{sec:impl} with the full scaling
sweep, the per-phase decomposition, and the head-to-head with Binius64 that
substantiates the comparison of \Cref{sec:cmp}. All figures use one Apple~M4
($10$-core, $16$\,GB), CPU, at the matched commitment rate $1/4$; absolute times
are thermally variable ($\sim$$10$--$20\%$ run-to-run), so \emph{ratios} and
per-phase \emph{shapes} are the stable signal.

\subsection{The construction's scaling}
\label{sec:app-scaling}

\Cref{tab:zinc-scaling} sweeps the deployed prover, verifier, and proof across
trace heights at the proof-size-optimal matrix shape
($\mathrm{num\_rows}\approx\sqrt N$). The prover is near-linear and the verifier
$\approx\sqrt N$ until $\nu\ge 23$, where the working set exceeds the $16$\,GB
host and both turn memory-bound; the compute stays $O(N)$, so on a machine with
adequate RAM the linear regime extends. Proof size grows as $O(\sqrt N)$: the
balanced shape trades the combined row against the $t$ column openings, each
within $\sqrt2$ of the other, so the proof is not flat --- from $0.55$ MiB at one
block to $\approx 40$ MiB at $2^{24}$ rows.

\begin{table}[h]
\centering
\renewcommand{\arraystretch}{1.15}
\begin{tabular}{r r r r r r}
\toprule
$\nu$ & $\approx$comp. & $\mathrm{num\_rows}$ & prove & verify & proof \\
\midrule
$9$  & $7$         & $2$   & $2.2$ ms  & $1.2$ ms  & $0.55$ MiB \\
$14$ & $241$       & $16$  & $13.6$ ms & $3.8$ ms  & $1.62$ MiB \\
$16$ & $964$       & $32$  & $46.6$ ms & $8.7$ ms  & $2.88$ MiB \\
$19$ & $7{,}710$   & $64$  & $354$ ms  & $38$ ms   & $7.41$ MiB \\
$20$ & $15{,}420$  & $128$ & $667$ ms  & $64$ ms   & $10.33$ MiB \\
$21$ & $30{,}841$  & $128$ & $1.35$ s  & $117$ ms  & $14.36$ MiB \\
$22$ & $61{,}683$  & $256$ & $2.72$ s  & $224$ ms  & $20.21$ MiB \\
$23$ & $123{,}366$ & $256$ & $7.0$ s$^\dagger$  & $383$ ms$^\dagger$ & $28.23$ MiB \\
$24$ & $246{,}733$ & $512$ & $25.1$ s$^\dagger$ & $1.00$ s$^\dagger$ & $39.93$ MiB \\
\bottomrule
\end{tabular}
\caption{Deployed scaling: CPU, rate $1/4$, proof-size-optimal shape. $^\dagger$At
$\nu\ge 23$ the working set exceeds $16$\,GB, so prove/verify are memory-bound
(super-linear); proof size is unaffected. The prover nonetheless \emph{completes}
$123{,}366$ compressions (a $\approx$$17$\,GB working set, via swap), where
Binius64's $O(N)$ circuit build walls near $16{,}000$ (\Cref{tab:scaling}).}
\label{tab:zinc-scaling}
\end{table}

\subsection{Per-phase decomposition}
\label{sec:app-perphase}

\Cref{tab:perphase} decomposes the prover by phase. Every phase scales
near-linearly --- each $\approx 4\times$ per $4\times$ compressions, the total
tracking $N^{0.97}$ --- so there is no super-linear term, the property that lets
the prover run past the size where a non-uniform circuit's $O(N)$ build stalls.
The \emph{commit} scales steepest (its share grows $7\%\to30\%$, the large-$N$
centre of gravity and the phase the GPU commit most relieves); the \emph{open}
scales shallowest (its $t = 987$ openings are a fixed count, share shrinking
$17\%\to9\%$).

\begin{table}[h]
\centering
\renewcommand{\arraystretch}{1.2}
\begin{tabular}{l r r r c}
\toprule
Phase & $\nu{=}10$ & $\nu{=}16$ & $\nu{=}20$ & per $4\times$ \\
      & ($15$c) & ($963$c) & ($15{,}420$c) & comp \\
\midrule
Commit (encode $+$ Merkle)          & $0.4$ & $18.0$ & $337.9$ & $4.3\times$ \\
Linear PIOP (IC, $\psia$, sumcheck) & $1.9$ & $17.8$ & $255.3$ & $3.8\times$ \\
Hadamard discharge                  & $1.2$ & $19.0$ & $287.6$ & $3.9\times$ \\
Multipoint eval                     & $1.1$ & $11.2$ & $146.0$ & $3.6\times$ \\
Open                                & $0.9$ & $9.6$  & $102.7$ & $3.3\times$ \\
\midrule
Total                               & $5.3$ & $75.6$ & $1129$  & $3.9\times$ \\
\bottomrule
\end{tabular}
\caption{Prover by phase (ms), CPU, $t = 987$, across trace height. Near-linear
throughout ($4\times$ per $4\times$ compressions is exactly linear); commit
scales steepest, open shallowest. Per-phase \emph{shares} are stable across
thermal state.}
\label{tab:perphase}
\end{table}

\subsection{Head-to-head with Binius64}
\label{sec:app-binius}

\Cref{tab:scaling} sweeps both systems at the matched rate $1/4$ --- Binius64 by
compression count, the construction by $\nu$. Two divergences stand out. The
prover gap \emph{widens} from $\sim$$1.1\times$ at $10^3$ compressions to
$\sim$$2\times$ at $10^4$, and Binius64's example walls above $\sim$$8000$
(its $O(N)$ circuit build exceeds a minute). The verifiers split qualitatively:
Binius64's grows with $N$ (its public-input/wiring check is $O(N)$, reaching
$\sim$$0.5$\,s at $8000$ compressions) while the construction's grows like
$\sqrt N$ and stays in the tens of milliseconds.

\begin{table}[h]
\centering
\renewcommand{\arraystretch}{1.2}
\begin{tabular}{r r r r r}
\toprule
 & \multicolumn{2}{c}{Binius64 (rate $1/4$)} & \multicolumn{2}{c}{This work (CPU)} \\
\cmidrule(lr){2-3}\cmidrule(lr){4-5}
$\approx$compressions & prove & verify & prove & verify \\
\midrule
$\sim$$60$       & $33.7$ ms & $3.1$ ms  & $10.1$ ms & $4.8$ ms \\
$\sim$$250$      & $41.6$ ms & $4.6$ ms  & $23.3$ ms & $9.3$ ms \\
$\sim$$1000$     & $84.1$ ms & $14.0$ ms & $75.0$ ms & $17.9$ ms \\
$\sim$$2000$     & $155$ ms  & $27.5$ ms & $148$ ms  & $19.8$ ms \\
$\sim$$8000$     & $1.22$ s  & $493$ ms  & $581$ ms  & $45$ ms \\
$\sim$$15{,}400$ & \multicolumn{2}{c}{(build $>1$ min)} & $1.19$ s & $83$ ms \\
\bottomrule
\end{tabular}
\caption{SHA-256 scaling at rate $1/4$, CPU, representative single runs.
\emph{Proof sizes}: Binius64 $123\!\to\!299$ KiB ($16\!\to\!8192$ compressions;
recursive FRI grows slowly); this work $O(\sqrt N)$ --- $0.55/2.88/10.3$ MiB at
$\approx 7/964/15{,}420$ compressions. The $O(N)$-vs-$\sqrt N$ verifier split is
the robust signal.}
\label{tab:scaling}
\end{table}

The mechanism is in the per-phase decomposition. \Cref{tab:perphase-b64} shows
Binius64's prover is \emph{super-linear} and worsening ($\sim$$4.6\times$ per
$4\times$ compressions around $10^3$, past $6\times$ by $8000$, versus this
work's flat $\sim$$3.9\times$): the BitAnd and Shift reductions build
$\mathrm{GF}(2^{128})$ tensors whose working set spills cache. The Shift reduction
dominates at small $N$ but its share falls ($62\%\to32\%$); the BitAnd reduction
rises to dominate ($14\%\to41\%$) and drives the super-linear growth.

\begin{table}[h]
\centering
\renewcommand{\arraystretch}{1.2}
\begin{tabular}{l r r r r}
\toprule
Phase & $64$c & $1024$c & $4096$c & $8192$c \\
\midrule
Commit (FRI)        & $1.8$  & $15.9$ & $86.3$  & $144.7$ \\
BitAnd reduction    & $4.0$  & $32.7$ & $219.4$ & $557.1$ \\
Shift reduction     & $16.5$ & $51.9$ & $171.2$ & $434.7$ \\
PCS opening         & $3.9$  & $13.4$ & $45.9$  & $212.7$ \\
\midrule
Total prove         & $27.3$ & $116$  & $532$   & $1370$ \\
\bottomrule
\end{tabular}
\caption{Binius64 prover by phase (ms), rate $1/4$. The BitAnd reduction
overtakes the Shift reduction at large $N$ and drives super-linear growth.}
\label{tab:perphase-b64}
\end{table}

\paragraph{Preprocessing.} The starkest structural difference is not in any timed
phase but in the \emph{preprocessing} (\Cref{tab:intro-preproc}): Binius64's
non-uniform word circuit carries an $O(N)$ constraint-plus-wiring description
reaching hundreds of megabytes, its $\ge 16{,}384$-compression build exceeding a
minute, whereas this work's uniform AIR has an $O(1)$, kilobyte-scale description
independent of trace height. This $O(N)$-vs-$O(1)$ gap is the root of both
Binius64's $O(N)$ verifier and its circuit-build wall.
